\section{Sets as membership predicates}\label{ch16:sec:predicates}
A \emph{predicate} on a set $A$ is a rule that turns each element of $A$ into a proposition. For example, ``$n$ is even'' is a predicate on the integers. For each particular integer it gives a statement, such as ``$6$ is even'' or ``$7$ is even,'' and each of these statements is either true or false. You have already met a predicate in Lean: in Section~\ref{ch15:sec:exists}, the input to \code{Exists} was the function \code{fun n : Nat => n + 1 = 4}, which turns each natural number $n$ into the proposition $n+1=4$. In Lean, a predicate on a type \code{A} is a function of type \code{A -> Prop}.

A predicate describes a set, namely the set of elements for which it holds; in the set-builder notation of Section~\ref{ch03:sec:sets}, a predicate $P$ on $A$ describes the set $\{x\in A:P(x)\}$. Conversely, a set $S\subseteq A$ gives the predicate ``$x\in S$.'' So a subset of a type \code{A} can be represented in Lean by a predicate \code{s : A -> Prop}, where the proposition \code{s x} means that $x$ belongs to the set. Chapter~\ref{ch:maps} described a set by its members, and here we describe it by its membership condition. The two descriptions carry the same information, because the question ``does $x$ belong?'' has the same answer in both.

A predicate states what membership means. It does not have to come with a quick way of deciding membership. For example, ``some sign sequence of length $n$ has all of its periodic correlations at nonzero shifts equal to zero'' is a perfectly good predicate on the positive integers $n$, yet deciding exactly which $n$ satisfy it is the research problem of Section~\ref{ch14:sec:research}.

Here are the intersection and the preimage, defined through their membership conditions, together with a theorem that connects them:
\par\noindent\begin{minipage}{\linewidth}
\begin{lstlisting}
def meet {A : Type} (s t : A -> Prop) : A -> Prop :=
  fun x => And (s x) (t x)

def preimage {A B : Type} (f : A -> B)
    (s : B -> Prop) : A -> Prop :=
  fun x => s (f x)

theorem preimage_meet {A B : Type} (f : A -> B)
    (s t : B -> Prop) (x : A) :
    Iff (preimage f (meet s t) x)
      (meet (preimage f s) (preimage f t) x) := by
  rfl
\end{lstlisting}
\end{minipage}\par

Suppose \code{s} and \code{t} describe the sets $S$ and $T$. The predicate \code{meet s t} holds at \code{x} when both \code{s x} and \code{t x} hold, so it describes the intersection $S\cap T$. The predicate \code{preimage f s} holds at \code{x} when \code{s (f x)} holds, that is, when $f(x)\in S$, so it describes the preimage $f^{-1}(S)=\{x\in A:f(x)\in S\}$ of Chapter~\ref{ch:maps}. The expression \code{Iff P Q} is Lean's way of writing $P\Leftrightarrow Q$, ``$P$ if and only if $Q$.'' So the theorem \code{preimage\_meet} is the identity
\[
f^{-1}(S\cap T)=f^{-1}(S)\cap f^{-1}(T)
\]
from Chapter~\ref{ch:maps}, stated one element at a time: for each \code{x}, membership in the left side is equivalent to membership in the right side. Since two sets are equal exactly when they have the same members, this is the same as the equality of the two sets.

Why does \code{rfl} prove it? Unfold the definitions on the left side: \code{preimage f (meet s t) x} becomes \code{meet s t (f x)}, which becomes \code{And (s (f x)) (t (f x))}. Unfold the right side: \code{meet (preimage f s) (preimage f t) x} becomes \code{And (preimage f s x) (preimage f t x)}, which becomes the very same \code{And (s (f x)) (t (f x))}. Section~\ref{ch15:sec:symbols} used \code{rfl} to prove an equality whose two sides reduce to the same expression. The same tactic also proves an equivalence whose two sides reduce to the same proposition, and Lean unfolds \code{preimage} and \code{meet} by itself while it compares them. No rewriting is needed.

The proof is short because the definitions already contain the reasoning. In Chapter~\ref{ch:maps}, every step of the chain of equivalences for $f^{-1}(S\cap T)$ was justified by ``definition of preimage'' or ``definition of intersection.'' Those are exactly the unfoldings that Lean carries out, so the single line \code{rfl} covers the whole chain.

\pauseq{Suppose the right side of \code{preimage\_meet} were written \code{meet (preimage f t) (preimage f s) x}, with \code{s} and \code{t} exchanged. The statement is still true. Would \code{rfl} still prove it?}{No. The left side still unfolds to \code{And (s (f x)) (t (f x))}, but the right side now unfolds to \code{And (t (f x)) (s (f x))}. These two propositions are equivalent, but they are not the same expression, so Lean reports that \code{rfl} failed. A proof would have to swap the two halves of the conjunction, as the theorem \code{and\_swap} of Section~\ref{ch15:sec:conjunction} does.}

This example also shows why definitions deserve close attention. Suppose \code{meet} had been defined with \code{Or} in place of \code{And}. The theorem \code{preimage\_meet} would keep its name, its proof \code{rfl} would still work, and Lean would still accept it. But it would now be the statement about unions from Exercise~\ref{ex:16.2}, and the name \code{meet} would no longer describe what the definition does. A name is only a label. What a statement says is determined by the bodies of its definitions.
